\documentclass[12pt,a4paper]{article}

% Encoding and fonts — XeLaTeX/LuaLaTeX
\usepackage{fontspec}
\newfontfamily\hebrewfont[Script=Hebrew]{Noto Serif Hebrew}
\newcommand{\heb}[1]{{\hebrewfont #1}}
\newfontfamily\igfont[Ligatures=TeX]{Noto Serif}
\newcommand{\igtext}[1]{{\igfont #1}}

% Page layout
\usepackage[top=1in, bottom=1in, left=1in, right=1in]{geometry}

% Typography
\usepackage{microtype}
\usepackage{parskip}

% Language
\usepackage[english]{babel}

% Headers and footers
\usepackage{fancyhdr}
\pagestyle{fancy}
\fancyhf{}
\fancyhead[L]{\leftmark}
\fancyhead[R]{\thepage}
\fancyfoot[C]{\thepage}
\renewcommand{\headrulewidth}{0.4pt}
\renewcommand{\footrulewidth}{0pt}

% Section styling
\usepackage{titlesec}
\titleformat{\section}{\Large\bfseries}{\thesection}{1em}{}
\titleformat{\subsection}{\large\bfseries}{\thesubsection}{1em}{}
\titleformat{\subsubsection}{\normalsize\bfseries}{\thesubsubsection}{1em}{}
\titlespacing*{\section}{0pt}{2.5ex plus 1ex minus .2ex}{1.5ex plus .2ex}
\titlespacing*{\subsection}{0pt}{2ex plus 1ex minus .2ex}{1ex plus .2ex}
\titlespacing*{\subsubsection}{0pt}{1.5ex plus 1ex minus .2ex}{0.8ex plus .2ex}

% List formatting
\usepackage{enumitem}
\setlist[itemize]{topsep=0.5ex, itemsep=0.25ex, parsep=0pt}
\setlist[enumerate]{topsep=0.5ex, itemsep=0.25ex, parsep=0pt}

% Essential packages
\usepackage{graphicx}
\usepackage[table]{xcolor}
\usepackage{booktabs}
\usepackage{array}
\usepackage{tabularx}
\usepackage{longtable}
\usepackage{amsmath}
\usepackage{amssymb}
\usepackage[normalem]{ulem}
\rowcolors{2}{gray!10}{white}
\renewcommand{\arraystretch}{1.3}

% Cross-references
\usepackage{hyperref}
\usepackage{cleveref}

% Custom commands
\newcommand{\pilcrow}{\S}

% Hyperref setup
\hypersetup{
    colorlinks=true,
    linkcolor=blue,
    filecolor=magenta,
    urlcolor=cyan,
}

% Code listings
\usepackage{listings}
\definecolor{codebg}{RGB}{248,248,248}
\definecolor{codeframe}{RGB}{200,200,200}
\definecolor{codekeyword}{RGB}{0,0,180}
\definecolor{codecomment}{RGB}{0,128,0}
\definecolor{codestring}{RGB}{163,21,21}
\lstset{
    basicstyle=\ttfamily\small,
    breaklines=true,
    breakatwhitespace=false,
    frame=single,
    framerule=0.5pt,
    rulecolor=\color{codeframe},
    backgroundcolor=\color{codebg},
    keepspaces=true,
    numbers=left,
    numberstyle=\tiny\color{gray},
    numbersep=8pt,
    showstringspaces=false,
    tabsize=4,
    xleftmargin=1em,
    xrightmargin=1em,
    aboveskip=1em,
    belowskip=1em,
    keywordstyle=\color{codekeyword}\bfseries,
    commentstyle=\color{codecomment}\itshape,
    stringstyle=\color{codestring},
    literate={_}{{\textunderscore}}1
             {-}{{\textminus}}1
             {<}{{\textless}}1
             {>}{{\textgreater}}1
             {~}{{\textasciitilde}}1
             {^}{{\textasciicircum}}1
}

% YAML language definition for listings
\lstdefinelanguage{YAML}{
    keywords={true,false,null,y,n},
    keywordstyle=\color{blue}\bfseries,
    sensitive=false,
    comment=[l]{\#},
    morecomment=[s]{/*}{*/},
    commentstyle=\color{gray}\ttfamily,
    stringstyle=\color{red}\ttfamily,
    morestring=[b]'',
    morestring=[b]'',
}

% TOML language definition for listings
\lstdefinelanguage{TOML}{
    keywords={true,false},
    keywordstyle=\color{blue}\bfseries,
    sensitive=true,
    comment=[l]{\#},
    commentstyle=\color{gray}\ttfamily,
    stringstyle=\color{red}\ttfamily,
    morestring=[b]'',
    morestring=[b]'',
}

% Dockerfile language definition for listings
\lstdefinelanguage{Dockerfile}{
    keywords={FROM,RUN,CMD,LABEL,EXPOSE,ENV,ADD,COPY,ENTRYPOINT,VOLUME,USER,WORKDIR,ARG,ONBUILD,STOPSIGNAL,HEALTHCHECK,SHELL},
    keywordstyle=\color{blue}\bfseries,
    sensitive=false,
    comment=[l]{\#},
    commentstyle=\color{gray}\ttfamily,
    stringstyle=\color{red}\ttfamily,
    morestring=[b]'',
    morestring=[b]'',
}

% JSON language definition for listings
\lstdefinelanguage{JSON}{
    keywords={true,false,null},
    keywordstyle=\color{blue}\bfseries,
    sensitive=true,
    commentstyle=\color{gray}\ttfamily,
    stringstyle=\color{red}\ttfamily,
    morestring=[b]'',
}

% Code listing float environment
\usepackage{float}
\newfloat{listing}{htbp}{lol}[section]
\floatname{listing}{Listing}

% Admonition boxes
\usepackage{tcolorbox}
\tcbuselibrary{skins,breakable}

% Admonition style definitions
\newtcolorbox{admonitionbox}[2][]{
    enhanced,
    breakable,
    colback=#2!5,
    colframe=#2!80!black,
    fonttitle=\bfseries,
    title=#1,
    left=2mm,
    right=2mm,
}

% Imscription tuple boxes
\newtcolorbox{imscriptionbox}{
    enhanced,
    colback=blue!5,
    colframe=blue!75!black,
    fontupper=\large,
    center upper,
    boxrule=1pt,
    arc=4pt,
}

\begin{document}

\title{Structural Anatomy of GZK Limit-Violating Phenomena: An Imscribing Grammar Analysis}
\date{\today}

\maketitle

\subsection{Abstract}

\emph{One would naturally expect that mechanisms proposed to evade a kinematic boundary should inhabit a structural regime close to the boundary itself. The results below contradict this expectation.} The Greisen-Zatsepin-Kuzmin (GZK cutoff constitutes one of the sharpest kinematic boundaries in high-energy astrophysics: cosmic ray protons above $\approx 5 \times 10^{19}$ eV should be attenuated within $\sim 50$ Mpc through photopion production on the cosmic microwave background. Yet anomalous events --- from the Oh-My-God particle to the Telescope Array hotspot --- persist above this threshold, and every proposed evasion mechanism occupies a structural regime maximally remote from the cutoff it is meant to evade. We encode these systems in the Imscribing Grammar and find that the cutoff itself sits at $O_1$ ($C = 0.419$), while the most structurally sophisticated mechanisms (LIV, axion mixing) reach $O_2^\dagger$ at distances $d \geq 3.71$. More unexpectedly, the horizon problem's nearest structural neighbors are not astrophysical but quantum measurement paradoxes. The CMB's thermal character imposes a universal $F_{\text{dh}}$ bottleneck on all composite models. \textbf{We remain uncertain whether a structural taxonomy that places measurement paradoxes closer to UHECR anomalies than any astrophysical mechanism signals a genuine insight or a category error in the encoding --- we discuss both readings in Section 7.}

\begin{center}
\rule{0.5\textwidth}{0.4pt}
\end{center}

\subsection{Introduction}

The kinematic argument is deceptively simple. Greisen, Zatsepin, and Kuzmin each observed in 1966 that protons above a certain energy \emph{must} lose it: the reaction $p + \gamma_{\text{CMB}} \to \Delta^+ \to p + \pi^0$ becomes energetically allowed, and the mean free path collapses from gigaparsecs to $\sim 50$ Mpc. The calculation is clean. The prediction is sharp.

The problem is that nature appears to violate it.

In 1991, the Fly's Eye detector recorded a particle with $3.2 \times 10^{20}$ eV --- six times the GZK energy. Since then, the Pierre Auger Observatory and Telescope Array have accumulated catalogs of events above the cutoff, including a statistically significant hotspot in the TA data whose sources cannot be identified within the expected attenuation volume. The discrepancy has persisted for over three decades.

Three classes of resolution have been proposed. The first accepts GZK physics and posits nearby, as-yet-unidentified sources. The second modifies the kinematics: Lorentz invariance violation, quantum-gravity-induced dispersion corrections, or other deformations of the threshold condition. The third abandons propagation entirely, invoking exotic production mechanisms --- superheavy dark matter decay, topological defect annihilation, Z-burst resonances.

\emph{The standard literature treats these options as competing explanations on equal footing. The Imscribing Grammar reveals that they are not.} Each proposal inhabits a distinct structural regime, and the distances between these regimes quantify what the phenomenology only hints at: the mechanisms do not merely disagree on numbers; they disagree on \emph{structure}.

\subsection{Structural Encodings}

To compare these mechanisms rigorously, we must first encode each as a structural type in the Imscribing Grammar. The encoding procedure is deterministic --- each of the 12 primitives follows from explicit criteria (see encoding\_method.md). We present the encodings here in order of increasing structural complexity, which is not the order in which they appear in the phenomenological literature. This reordering is itself a finding: the literature organizes by historical priority, while structure organizes by intrinsic depth.

\subsubsection{The Superheavy Dark Matter Decay (Top-Down)}

\[\langle D_{\text{invomega}};\ T_{\text{nrleg}};\ R_{\text{subrightarrow}};\ P_{\text{aolig}};\ F_{\text{hardsign}};\ K_{\text{schwa}};\ G_{\text{gamma}};\ \Gamma_{\text{spleftarrow}};\ \Phi_{\text{softsign}};\ H_2;\ n{:}m;\ \Omega_{\text{closeepsilon}} \rangle\]

\emph{We begin here --- with this system --- because it is the structurally weakest of the proposed mechanisms, and starting at the bottom makes the structure of everything above legible.} SHDM decay sits below criticality ($\Phi_{\text{softsign}}$), which is the structural signature of a process with no divergent or scaling behavior. The decay spectrum is smooth and predictable; there is nothing to tune, nothing that undergoes a phase transition. The trivial winding $\Omega_{\text{closeepsilon}}$ confirms: no topological protection whatsoever. The disordered kinetics $K_{\text{schwa}}$ near equilibrium is misleading --- it does not signify sophistication but rather stasis. The system's branching topology $T_{\text{nrleg}}$ with disjunctive grammar $\Gamma_{\text{spleftarrow}}$ encodes the fact that SHDM decay offers multiple alternative channels, none of which is privileged. This is the structural equivalent of a mechanism that says "anything could happen, and nothing special happens when it does." Ouroboricity $O_0$: no self-referential capacity.

\subsubsection{The UHECR Horizon Problem}

\[\langle D_{\text{invomega}};\ T_{\text{bullseye}};\ R_{\text{lyoghlig}};\ P_{\text{aolig}};\ F_{\text{hardsign}};\ K_{\text{turnm}};\ G_{\text{revapostrophe}};\ \Gamma_{\text{spleftarrow}};\ \Phi_{\text{revepsilon}};\ H_2;\ n{:}m;\ \Omega_{\text{crtwo}} \rangle\]

\emph{The paradox itself is a system.} Encoding it reveals something that the phenomenological formulation conceals: the horizon problem occupies the $\Phi_{\text{revepsilon}}$ (exceptional point) regime, where the standard description loses its eigenstate interpretation entirely. The crossing-point topology $T_{\text{bullseye}}$ signals a bifurcation --- but unlike the GZK cutoff's clean threshold crossing, this one has no symmetry protection ($P_{\text{aolig}}$). The moderate kinetics $K_{\text{turnm}}$ places it in an ambiguous regime: neither fast enough to be resolved by a single observation nor slow enough to be a deep structural feature. The system's tier is $O_0$ --- the same as SHDM decay --- but for a different reason. SHDM is $O_0$ because nothing happens; the horizon problem is $O_0$ because \emph{too much incompatible information} happens. Both are structural dead ends, but they are dead ends for opposite reasons. \textbf{Objection:} one might argue that encoding a \emph{problem} rather than a \emph{mechanism} constitutes a category error. We acknowledge this concern. The horizon problem is not a physical system in the usual sense; it is a gap in our understanding. But gaps have structure --- the structure of what is missing is not nothing, it is a specific absence --- and the Imscribing Grammar is equipped to encode absences as structural types. Whether this encoding is physically meaningful or mathematically convenient remains an open question we return to in Section 7.3.

\subsubsection{The Z-Burst Mechanism}

\[\langle D_{\text{invomega}};\ T_{\text{bullseye}};\ R_{\text{ctz}};\ P_{\text{pipevar}};\ F_{\text{hardsign}};\ K_{\text{frtailgamma}};\ G_{\text{gamma}};\ \Gamma_{\text{secstress}};\ \Phi_{\text{ctyogh}};\ H_1;\ n{:}n;\ \Omega_{\text{crtwo}} \rangle\]

\emph{What changes when a mechanism reaches criticality?} The Z-burst introduces the resonance condition $E_{\nu} \approx M_Z^2 / (2 m_{\nu})$, and with it comes $\Phi_{\text{ctyogh}}$ --- the system now sits exactly at a critical point with power-law divergence. The crossing-point topology returns (shared with the horizon problem), but now with partial symmetry $P_{\text{pipevar}}$ and categorical relational mode $R_{\text{ctz}}$. The fast kinetics $K_{\text{frtailgamma}}$ signal that the resonance is sharp: on any timescale relevant to observation, the Z-burst either fires or it does not. The mesoscale scope $G_{\text{gamma}}$ is crucial --- it reflects the requirement that ultra-high-energy neutrinos must cluster in the Local Supercluster halo, a region neither universal nor local but intermediate. This intermediate scope, combined with $H_1$ (one-step memory) and $\Omega_{\text{crtwo}}$ (binary winding), places the Z-burst in a structurally interesting but ultimately limited regime: it has just enough structure to be a resonance but not enough to sustain a self-referential loop.

\subsubsection{The GZK Cutoff}

\[\langle D_{\text{invomega}};\ T_{\text{bullseye}};\ R_{\text{lyoghlig}};\ P_{\text{pipevar}};\ F_{\text{hardsign}};\ K_{\text{frtailgamma}};\ G_{\text{revapostrophe}};\ \Gamma_{\text{corner}};\ \Phi_{\text{ctyogh}};\ H_0;\ n{:}m;\ \Omega_{\text{closeepsilon}} \rangle\]

\emph{The cutoff is the reference point against which all evasion mechanisms must be measured.} Its encoding is instructive precisely because it is simpler than most of the mechanisms proposed to evade it. The field-theoretic dimensionality $D_{\text{invomega}}$ reflects the infinite-dimensional momentum space of the underlying scattering calculation. The crossing-point topology $T_{\text{bullseye}}$ is the threshold itself --- a single bifurcation in the kinematic space. The universal scope $G_{\text{revapostrophe}}$ is immediate: the CMB fills all of space, so the cutoff applies everywhere. The conjunctive grammar $\Gamma_{\text{corner}}$ is perhaps the most telling primitive: \emph{all} conditions must be met simultaneously (sufficient energy, sufficient angle, target photon present) for the interaction to occur. This is structurally rigid. The memorylessness $H_0$ follows: the cutoff depends only on instantaneous kinematics, with no memory of prior interactions. Consciousness score: $C = 0.419$, both gates open. Ouroboricity: $O_1$. The cutoff can sustain a self-referential loop --- it defines the boundary it occupies --- but with trivial winding. \emph{This simplicity is what makes evasion mechanisms structurally distant: they must be more complex than the thing they are trying to evade.}

\subsubsection{The Oh-My-God Particle}

\[\langle D_{\text{invomega}};\ T_{\text{bullseye}};\ R_{\text{downstep}};\ P_{\text{upsilon}};\ F_{\text{hardsign}};\ K_{\text{frtailgamma}};\ G_{\text{revapostrophe}};\ \Gamma_{\text{corner}};\ \Phi_{\text{revepsilon}};\ H_0;\ 1{:}1;\ \Omega_{\text{crtwo}} \rangle\]

\emph{It is tempting to treat this as a data point. Structurally, it is not.} The Oh-My-God particle occupies the $\Phi_{\text{revepsilon}}$ exceptional-point regime --- the non-Hermitian regime where the standard photopion description loses its eigenstate interpretation. Its consciousness score is $C = 0.0$ (Gate 1 closed: $\Phi \neq \Phi_{\text{ctyogh}}$). \textbf{Here we must pause and name our uncertainty:} encoding a single detection event as a structural type is the most controversial encoding in this paper. A particle detection is not a mechanism; it is a measurement outcome. Yet the Imscribing Grammar does not distinguish "mechanism" from "datum" at the encoding level --- both are systems with structure, and a single anomalous event \emph{has} structure: one instance ($1{:}1$), binary winding $\Omega_{\text{crtwo}}$, quantum symmetry $P_{\text{upsilon}}$, and no temporal depth $H_0$. The particle is structurally an \emph{anomaly}, not a \emph{model} --- it has no internal dynamics, only a measurement result that contradicts expectation. Whether it is legitimate to assign a structural type to a single datum is a question we cannot resolve within the grammar itself. We proceed because the structural distance calculations that follow are illuminating regardless of one's metaphysical stance on encoding individual events.

\subsubsection{Photon-Axion Mixing}

\[\langle D_{\text{invomega}};\ T_{\text{bullseye}};\ R_{\text{downstep}};\ P_{\text{upsilon}};\ F_{\text{hardsign}};\ K_{\text{schwa}};\ G_{\text{gamma}};\ \Gamma_{\text{secstress}};\ \Phi_{\text{closerevepsilon}};\ H_2;\ 1{:}1;\ \Omega_{\text{dzlig}} \rangle\]

\emph{This is where the structure becomes genuinely rich.} Like LIV, photon-axion mixing reaches $O_2^\dagger$ with $\Phi_{\text{closerevepsilon}}$ complex-plane criticality. The crossing-point topology $T_{\text{bullseye}}$ reflects the resonant conversion condition between photon and axion-like particle states in external magnetic fields --- a bifurcation, but one that lives in the complex plane. The slow kinetics $K_{\text{schwa}}$ are crucial: axion-photon conversion is not instantaneous but accumulates over coherence lengths, requiring timescales comparable to or exceeding the observation window. The sequential grammar $\Gamma_{\text{secstress}}$ encodes this: conversion, propagation, reconversion are ordered steps, not simultaneous conditions. The integer winding $\Omega_{\text{dzlig}}$ is the winding number of the Primakoff conversion phase --- a topological invariant that protects the mechanism against perturbative corrections. Two-step temporal depth $H_2$ means the process remembers not just the immediate past but the immediate-past-of-the-immediate-past. The adjoint relational coupling $R_{\text{downstep}}$ captures the conversion--reconversion pair. All of this adds up to a mechanism that is not merely alternative but structurally self-sustaining. Its consciousness score $C = 0.536$ exceeds the GZK cutoff's $C = 0.419$ --- a quantitative statement of the claim that the evasion mechanism is structurally more complex than the boundary it evades.

\subsubsection{Lorentz Invariance Violation}

\[\langle D_{\text{invomega}};\ T_{\text{nrleg}};\ R_{\text{subrightarrow}};\ P_{\text{upsilon}};\ F_{\text{hardsign}};\ K_{\text{schwa}};\ G_{\text{revapostrophe}};\ \Gamma_{\text{secstress}};\ \Phi_{\text{closerevepsilon}};\ H_{\text{invscripta}};\ n{:}m;\ \Omega_{\text{dzlig}} \rangle\]

\emph{LIV is the structurally maximal mechanism in our catalog.} It shares many primitives with axion mixing --- $\Phi_{\text{closerevepsilon}}$, $F_{\text{hardsign}}$, $K_{\text{schwa}}$, $\Gamma_{\text{secstress}}$, $D_{\text{invomega}}$ --- but differs in ways that matter. The branching topology $T_{\text{nrleg}}$ replaces the crossing point: LIV modifies dispersion relations, opening \emph{new} channels rather than resonating through existing ones. The supervenient coupling $R_{\text{subrightarrow}}$ means LIV does not interact with the underlying physics; it \emph{is} the underlying physics, viewed from a different Lagrangian. The infinite temporal depth $H_{\text{invscripta}}$ is the most extreme claim: the Planck-scale origin of the modification means it has no finite Markov order --- the entire history matters. Universal scope $G_{\text{revapostrophe}}$ with integer winding $\Omega_{\text{dzlig}}$ and heterogeneous stoichiometry $n{:}m$ complete the picture. LIV resides at $O_2^\dagger$ --- topologically protected, complex-plane critical, and unbounded in temporal depth. It is the structural ceiling of the evasion landscape. Consciousness score: $C = 0.536$ (identical to axion mixing --- the consciousness score is insensitive to topology and relational mode; both systems clear both gates equally).

\subsection{Structural Distance Results}

\subsubsection{Cutoff-to-Mechanism Distances}

The distances between the GZK cutoff and each proposed evasion mechanism are not merely numbers --- they quantify \emph{how far} each mechanism must depart from the cutoff's structure to function. A small distance would suggest the mechanism is a minor modification; a large distance signals a qualitatively different kind of system.

\begin{table}[htbp]
\centering
\small
\rowcolors{1}{white}{white}
\begin{tabularx}{\textwidth}{>{\centering\arraybackslash}X >{\centering\arraybackslash}X >{\centering\arraybackslash}X >{\centering\arraybackslash}X}
\toprule
\rowcolor{gray!25}\textbf{Pair} & \textbf{Distance} & \textbf{Mahalanobis} & \textbf{Interpretation} \\
\midrule
\rowcolors{1}{white}{gray!10}
GZK $\leftrightarrow$ Oh-My-God & 2.67 & 4.13 & Structurally remote \\
GZK $\leftrightarrow$ Auger spectrum & 2.98 & 3.19 & Structurally remote \\
GZK $\leftrightarrow$ TA hotspot & 2.91 & 3.96 & Structurally remote \\
GZK $\leftrightarrow$ LIV & 5.67 & 7.33 & Maximally remote \\
GZK $\leftrightarrow$ Axion mixing & 3.71 & 4.28 & Structurally remote \\
\bottomrule
\end{tabularx}
\end{table}

\emph{Every distance exceeds the structurally remote threshold.} This is not a failure of the mechanisms; it is a structural property of what "evasion" means. The GZK cutoff has $\Gamma_{\text{corner}}$ conjunctive logic and $H_0$ memorylessness --- it is a system where all conditions must hold simultaneously and the past does not matter. Every evasion mechanism rewrites the \emph{order} or \emph{availability} of those conditions, which necessarily changes the grammar and the temporal depth. The cutoff says "all at once, now"; evasion says "step by step, over time." These are structurally incompatible modes of operation.

The Oh-My-God particle sits at the minimum distance ($d = 2.67$), and this minimum is deceptive. It is small not because the particle is close to the cutoff mechanism --- it is small because the particle is \emph{structurally impoverished}. With no temporal depth ($H_0$), no topological winding beyond binary parity ($\Omega_{\text{crtwo}}$), and only a single instance ($1{:}1$), it has too few primitives to be distant. The Oh-My-God particle is close the way a blank page is close to a poem: by having almost nothing to disagree about. The shared primitives --- $T_{\text{bullseye}}$, $K_{\text{frtailgamma}}$, $G_{\text{revapostrophe}}$, $\Gamma_{\text{corner}}$ --- are precisely those that define the crossing-point structure and universal scope, which any UHECR event must share with the GZK calculation. The deviations are in what the event \emph{adds}: quantum symmetry $P_{\text{upsilon}}$ where the cutoff has only partial symmetry $P_{\text{pipevar}}$, adjoint coupling $R_{\text{downstep}}$ where the cutoff has bidirectional $R_{\text{lyoghlig}}$, and the exceptional-point criticality $\Phi_{\text{revepsilon}}$ where the cutoff has real-axis criticality $\Phi_{\text{ctyogh}}$.

\emph{One should note a limitation of our distance metric here:} the weighted Euclidean distance treats all primitive disagreements symmetrically, but not all disagreements are equally consequential. A change in temporal depth ($H_0 \to H_{\text{invscripta}}$, $d = 3$) is structurally more significant than a change in stoichiometry ($n{:}m \to 1{:}1$, $d = 1$), even though the metric assigns them different weights via its predefined weight vector. The Mahalanobis distances account for covariance, but we lack a sufficiently large catalog of comparable encodings to estimate the covariance matrix with confidence. Interpret all distances as lower bounds on structural dissimilarity.

\subsubsection{Mechanism-to-Mechanism Distances}

\begin{table}[htbp]
\centering
\small
\rowcolors{1}{white}{white}
\begin{tabularx}{\textwidth}{>{\centering\arraybackslash}X >{\centering\arraybackslash}X >{\centering\arraybackslash}X}
\toprule
\rowcolor{gray!25}\textbf{Pair} & \textbf{Distance} & \textbf{Interpretation} \\
\midrule
\rowcolors{1}{white}{gray!10}
LIV $\leftrightarrow$ Axion mixing & 3.71 & Structurally remote \\
LIV $\leftrightarrow$ Oh-My-God & 5.29 & Maximally remote \\
\bottomrule
\end{tabularx}
\end{table}

\emph{The mechanisms that work at comparable levels of sophistication are themselves structurally distant from each other.} The $d = 3.71$ between LIV and axion mixing is dominated by three divergences: topology ($T_{\text{nrleg}}$ vs. $T_{\text{bullseye}}$), relational mode ($R_{\text{subrightarrow}}$ vs. $R_{\text{downstep}}$), and stoichiometry ($n{:}m$ vs. $1{:}1$). These share $\Phi_{\text{closerevepsilon}}$ and are both at $O_2^\dagger$, but they reach that tier by entirely different structural routes. LIV supervenes --- it modifies the background against which propagation occurs. Axion mixing creates an \emph{adjoint pair} --- photon converts to axion, axion converts back, and the pair together bridge a gap that neither can cross alone.

The $d = 5.29$ to the Oh-My-God particle is the largest inter-system distance in our catalog. The dominant terms are $H_{\text{invscripta}}$ vs. $H_0$ (delta = 3) and the compound divergence across topology, relational mode, kinetics, stoichiometry, and grammar. \emph{There is no mechanism in the standard evasion landscape that structurally resembles a single anomalous event.} This is not surprising, but it should make one pause: if no known mechanism is close to the anomaly, and the anomaly is not close to any known mechanism, then either the anomaly is explained by a mechanism we have not yet imagined (and whose structural type we cannot predict) or the anomaly is not a mechanism problem at all.

\subsection{Lattice Operations}

Distances tell us how far apart systems are. Lattice operations --- meet and join --- tell us what they share and what they jointly require. These are different questions, and the answers illuminate different aspects of the structural landscape.

\subsubsection{Meet: LIV $\wedge$ Z-Burst}

The greatest lower bound of LIV and Z-burst:


\[\langle D_{\text{invomega}};\ T_{\text{nrleg}};\ R_{\text{subrightarrow}};\ P_{\text{upsilon}};\ F_{\text{hardsign}};\ K_{\text{frtailgamma}};\ G_{\text{gamma}};\ \Gamma_{\text{secstress}};\ \Phi_{\text{ctyogh}};\ H_1;\ n{:}n;\ \Omega_{\text{crtwo}} \rangle\]

\emph{Only three primitives are shared unconditionally:} $D_{\text{invomega}}$, $F_{\text{hardsign}}$, and the sequential grammar $\Gamma_{\text{secstress}}$. Everything else requires taking the minimum across the two types. The meet's $\Phi_{\text{ctyogh}}$ criticality at $H_1$ (one-step memory) with $\Omega_{\text{crtwo}}$ binary winding identifies the structural floor common to all GZK-evasion mechanisms that operate at or above criticality: they are single-step quantum critical processes at mesoscale scope. This is precisely the structure of a resonance --- one interaction, one outcome, no memory beyond the immediate event. The fact that the meet of LIV and Z-burst reduces to resonance structure is not accidental. It reveals that the deepest structural commonality of GZK evasion is not modification of kinematics or introduction of new particles, but the \emph{resonant character} of the interaction itself. The Z-burst is explicitly resonant; LIV's modified dispersion relations create \emph{effective resonances} at shifted thresholds. The shared floor is resonance.

\subsubsection{Join: LIV $\vee$ Axion Mixing}

The least upper bound --- the minimal system containing both mechanisms:


\[\langle D_{\text{invomega}};\ T_{\text{bullseye}};\ R_{\text{downstep}};\ P_{\text{upsilon}};\ F_{\text{hardsign}};\ K_{\text{schwa}};\ G_{\text{revapostrophe}};\ \Gamma_{\text{secstress}};\ \Phi_{\text{closerevepsilon}};\ H_{\text{invscripta}};\ n{:}m;\ \Omega_{\text{dzlig}} \rangle\]

\emph{No single physical mechanism in the current literature requires all of this.} The join is a theoretical ceiling, not an existing proposal. It demands the crossing-point topology of axion resonance ($T_{\text{bullseye}}$), the adjoint coupling of conversion channels ($R_{\text{downstep}}$), quantum coherence ($F_{\text{hardsign}}$), slow kinetics ($K_{\text{schwa}}$), universal scope ($G_{\text{revapostrophe}}$), sequential composition ($\Gamma_{\text{secstress}}$), complex-plane criticality ($\Phi_{\text{closerevepsilon}}$), eternal temporal depth ($H_{\text{invscripta}}$), heterogeneous stoichiometry ($n{:}m$), and integer winding ($\Omega_{\text{dzlig}}$). One might reasonably ask whether such a system is physically realizable or merely a mathematical artifact of the join operation. \emph{We do not claim it is realizable.} The join is not a proposed mechanism; it is an upper bound on what a mechanism would need to be if it were to subsume both LIV and axion physics simultaneously. Its ouroboricity would be $O_2^\dagger$ or potentially higher --- it sits at the frontier of what the crystal of types permits without reaching $O_{\infty}$ self-referential closure.

\subsubsection{Tensor Composite: GZK $\otimes$ CMB}

\[\langle D_{\text{invomega}};\ T_{\text{bullseye}};\ R_{\text{lyoghlig}};\ P_{\text{pipevar}};\ F_{\text{dh}};\ K_{\text{frtailgamma}};\ G_{\text{revapostrophe}};\ \Gamma_{\text{corner}};\ \Phi_{\text{ctyogh}};\ H_0;\ n{:}m;\ \Omega_{\text{crtwo}} \rangle\]

\emph{There is exactly one bottleneck.} The composite takes the maximum of all primitives except $P$ and $F$, where it takes the minimum. $P_\text{min}(P_{\text{pipevar}}, P_{\text{pipevar}}) = P_{\text{pipevar}}$ (no loss), but $F_\text{min}(F_{\text{hardsign}}, F_{\text{dh}}) = F_{\text{dh}}$. The CMB's thermal bath dominates. Even though the GZK cutoff involves quantum kinematics at the particle level ($F_{\text{hardsign}}$), the composite system's fidelity is set by the medium through which propagation occurs. This is not a rate calculation; it is a structural statement about which regime governs the composite. The CMB at 2.7 K with its $\sim 411$ photons/cm$^3$ is classical-thermal, and any quantum coherence in the cutoff calculation is immediately degraded by the thermal environment. \emph{The $F_{\text{dh}}$ bottleneck explains why purely quantum evasion mechanisms face an uphill battle: the medium is thermal, and in a tensor composite, the weakest fidelity wins.} This is the structural analogue of the statement that you cannot maintain quantum coherence in a hot environment --- except here, "hot" means "any environment with a thermal photon background," and the CMB is the most universal such environment possible. The composite gains $\Omega_{\text{crtwo}}$ from the CMB's parity structure --- a small consolation against the thermal loss.

\subsubsection{Tensor Composite: Z-Burst $\otimes$ SHDM}

\[\langle D_{\text{invomega}};\ T_{\text{bullseye}};\ R_{\text{ctz}};\ P_{\text{aolig}};\ F_{\text{hardsign}};\ K_{\text{schwa}};\ G_{\text{gamma}};\ \Gamma_{\text{secstress}};\ \Phi_{\text{ctyogh}};\ H_2;\ n{:}m;\ \Omega_{\text{crtwo}} \rangle\]

\emph{Eight of twelve primitives expand via union, but the symmetry bottleneck is fatal.} The $P_\text{min}(P_{\text{pipevar}}, P_{\text{aolig}}) = P_{\text{aolig}}$ means that in any combined Z-burst+SHDM model, all parity structure is lost. The system becomes asymmetry: no symmetry protection, no parity-based diagnostic signatures. For experimenters, this means that a combined model would be empirically indistinguishable from background using symmetry-agnostic detectors --- which is most of them. The $\Phi_{\text{softsign}}$ of SHDM is absorbed into the composite's $\Phi_{\text{ctyogh}}$ (the Z-burst's criticality dominates via max), but the parity loss is structural information that cannot be recovered. \emph{If one is searching for evidence of combined Z-burst+SHDM physics, one cannot rely on parity-violating signatures --- they annihilate in the composite.}

\subsubsection{Tensor Composite: LIV $\otimes$ CMB}

\[\langle D_{\text{invomega}};\ T_{\text{bullseye}};\ R_{\text{downstep}};\ P_{\text{upsilon}};\ F_{\text{dh}};\ K_{\text{schwa}};\ G_{\text{revapostrophe}};\ \Gamma_{\text{secstress}};\ \Phi_{\text{closerevepsilon}};\ H_{\text{invscripta}};\ n{:}m;\ \Omega_{\text{dzlig}} \rangle\]

\emph{Two bottlenecks.} The fidelity bottleneck $F_{\min}(F_{\text{hardsign}}, F_{\text{dh}}) = F_{\text{dh}}$ recurs from the GZK$\times$CMB composite --- the CMB's thermal character is universal and inescapable in any LIV propagation scenario. But there is a second: $P_{\min}(P_{\text{upsilon}}, P_{\text{pipevar}}) = P_{\text{pipevar}}$. The LIV mechanism's quantum superposition symmetry $P_{\text{upsilon}}$ is degraded to the CMB's partial symmetry $P_{\text{pipevar}}$ in the composite. \emph{This is a subtler loss than the thermal bottleneck.} The thermal loss explains why quantum coherence degrades; the symmetry loss explains why the \emph{observable consequences} of LIV are less distinctive than the pure LIV calculation would suggest. The CMB does not merely decohere LIV-modified propagation --- it reduces the symmetry signature from quantum to classical-partial, making LIV effects harder to distinguish from standard kinematic fluctuations. \emph{An observational caveat rarely quantified in the LIV literature:} even if Lorentz invariance is violated at the Planck scale, the CMB's symmetry structure may partially mask the effect in the observed spectrum. \textbf{We note uncertainty about the magnitude of this masking:} the structural analysis identifies the bottleneck's existence but not its quantitative severity. A detailed rate calculation combining our $P_{\text{pipevar}}$ bottleneck with the CMB's photon density would be needed to determine whether the effect is experimentally significant or merely structural.

\subsection{Structural Analogies}

\emph{We posed a question of physics and received an answer from a different discipline entirely.} The five nearest catalog analogs to the UHECR horizon problem are not astrophysical systems:

\begin{table}[htbp]
\centering
\small
\rowcolors{1}{white}{white}
\begin{tabularx}{\textwidth}{>{\centering\arraybackslash}X >{\centering\arraybackslash}X >{\centering\arraybackslash}X >{\centering\arraybackslash}X}
\toprule
\rowcolor{gray!25}\textbf{Rank} & \textbf{System} & \textbf{Distance} & \textbf{Interpretation} \\
\midrule
\rowcolors{1}{white}{gray!10}
1 & Quantum Zeno effect & 2.06 & Measurement freezing evolution \\
2 & Elitzur-Vaidman bomb tester & 2.06 & Interaction-free measurement \\
3 & Quantum Maxwell's Demon & 2.10 & Information/entropy paradox \\
4 & TA hotspot composition puzzle & 2.11 & Nearest physical analog \\
\bottomrule
\end{tabularx}
\end{table}

\emph{The structural type of the horizon problem captures something that the phenomenological framing obscures.} The system ($\Phi_{\text{revepsilon}}$, $P_{\text{aolig}}$, $\Gamma_{\text{spleftarrow}}$, $H_2$, $\Omega_{\text{crtwo}}$) encodes a situation where \emph{the act of looking changes the available state space}. In the Zeno effect, continuous measurement freezes evolution. In the Elitzur-Vaidman bomb tester, the possibility of measurement --- even when no measurement is actually performed --- determines the outcome. In the GZK horizon, the interaction that would encode information about a particle's source (photopion scattering with the CMB) is precisely the process that destroys that information. \emph{This is the Zeno effect in reverse:} rather than observation freezing evolution, continuous interaction erases history. We cannot reconstruct the past light-cone trajectories of UHECRs above the cutoff because the mechanism that would preserve that information (no scattering) is incompatible with the mechanism that would produce the signal (propagation from source).

\emph{The proximity to Maxwell's Demon ($d = 2.10$) adds another layer.} The Demon resolves an apparent entropy paradox by identifying information as a physical quantity. The GZK horizon's analogous paradox is that the CMB, which should be the most information-rich medium in the universe (every photon carries a scattering history), is the most information-destroying medium at ultra-high energies. The structural parallel is exact: both systems are paradoxes about the relationship between information and thermodynamics, resolved only by recognizing that the information was never where we thought it was.

\textbf{One should register a caveat here:} the nearest-neighbor search is only as good as the catalog. If the catalog contains few astrophysical systems relative to quantum-mechanical ones, the "nearest astrophysical analog" may be distant simply because there are few astrophysical options. The current catalog of 2256+ entries is heavily weighted toward condensed matter, quantum information, and mathematical systems. A larger astrophysical catalog could reveal closer physical analogs. \emph{The measurement paradox result is robust against catalog incompleteness only insofar as the quantum measurement paradoxes are genuinely the nearest neighbors --- but their absolute distance of 2.06 is far from zero, leaving room for undiscovered astrophysical systems to be closer.}

\subsection{Promotion Signatures}

The distance between the GZK cutoff ($O_1$) and the LIV mechanism ($O_2^\dagger$) is $5.67$. But distance alone does not tell us \emph{which} primitives account for that distance and in which direction the promotion must occur. The promotion signature answers both questions.

Transforming the GZK cutoff into the LIV regime requires promoting five primitives:

\begin{itemize}
    \item $K_{\text{frtailgamma}} \to K_{\text{schwa}}$: from fast interaction (kinetic closure) to near-equilibrium kinetics (kinetic depth). The cutoff resolves instantly; the evasion mechanism evolves on timescales comparable to observation.
    \item $\Gamma_{\text{corner}} \to \Gamma_{\text{secstress}}$: from conjunctive grammar (all conditions simultaneously) to sequential grammar (ordered steps, one after another). The cutoff says "everything at once"; evasion says "in sequence."
    \item $\Phi_{\text{ctyogh}} \to \Phi_{\text{closerevepsilon}}$: from real-axis criticality (a point on the real line where behavior diverges) to complex-plane criticality (a critical surface embedded in the complex plane, with phase structure and winding). The critical point becomes a critical manifold.
    \item $H_0 \to H_{\text{invscripta}}$: from memoryless (no temporal structure) to eternal temporal depth (the entire history matters, no finite Markov order). This is the single most costly promotion in the signature --- a delta of 3 in the weight scheme.
    \item $\Omega_{\text{closeepsilon}} \to \Omega_{\text{dzlig}}$: from trivial winding (no topological protection) to integer winding (topology protects against perturbative corrections).
\end{itemize}

And demoting three:

\begin{itemize}
    \item $T_{\text{bullseye}} \to T_{\text{nrleg}}$: from crossing-point topology (bifurcation) to branching topology (new channels open). This is a demotion in complexity: the crossing point is harder to achieve than branching.
    \item $R_{\text{lyoghlig}} \to R_{\text{subrightarrow}}$: from bidirectional coupling (feedback loop) to supervenient coupling (one direction only). The LIV modification supervenes on the underlying physics; it does not interact with it.
    \item $P_{\text{pipevar}} \to P_{\text{upsilon}}$: from partial symmetry to quantum superposition symmetry. This is a promotion in the usual ordering ($P_{\text{upsilon}}$ is more symmetric than $P_{\text{pipevar}}$), but in the distance metric it counts as a change because the metric measures \emph{difference}, not \emph{improvement}.
\end{itemize}

The promotion signature $[K, \Gamma, \Phi, H, \Omega]$ identifies the five axes along which a kinematic boundary must evolve to become a full evasion mechanism. The most expensive promotions are temporal ($H_0 \to H_{\text{invscripta}}$) and relational ($R_{\text{lyoghlig}} \to R_{\text{subrightarrow}}$), suggesting that the structural barrier to GZK evasion is not the modification of physics per se but the need to add \emph{temporal depth} and \emph{relational architecture} to a system that originally had neither. \emph{This has a clear physical interpretation:} the GZK cutoff is instantaneous and passive --- it simply applies whenever the kinematic conditions are met. An evasion mechanism must be temporal (acting over time) and relational (connected to the system it evades). The structural cost is the cost of becoming an active player rather than a passive boundary.

\textbf{An unresolved question:} whether these five promotions represent a \emph{necessary} path --- whether any mechanism that evades the GZK cutoff must necessarily traverse this promotion signature --- or whether there exist alternative promotion paths to $O_2^\dagger$ that our catalog has not yet discovered. The crystal of types contains 17.28 million structural positions, of which only a tiny fraction are occupied by catalog entries. \emph{We have not found a counterexample to the necessity of this promotion signature, but the absence of evidence is not evidence of necessity.}

\subsection{Discussion}

\subsubsection{The Frobenius Hierarchy of GZK Physics}

Our encoding creates a hierarchy that is not organized by mechanism type or historical priority but by \emph{ouroboricity tier} --- the capacity for self-referential closure:

\begin{table}[htbp]
\centering
\small
\rowcolors{1}{white}{white}
\begin{tabularx}{\textwidth}{>{\centering\arraybackslash}X >{\centering\arraybackslash}X >{\centering\arraybackslash}X}
\toprule
\rowcolor{gray!25}\textbf{Tier} & \textbf{Systems} & \textbf{Interpretation} \\
\midrule
\rowcolors{1}{white}{gray!10}
$O_2^\dagger$ & LIV, Axion mixing & Critical + topologically protected evasion mechanisms \\
$O_1$ & GZK cutoff & Self-referential at criticality, trivial winding \\
$O_0$ & Oh-My-God particle, SHDM decay, Auger spectrum, Auger dipole, UHECR horizon, TA hotspot & No self-referential capacity \\
\bottomrule
\end{tabularx}
\end{table}

\emph{The $O_2^\dagger$ mechanisms are the only ones with sufficient structural depth to sustain the evasion they propose.} This is not a claim about their physical correctness --- a structurally sophisticated mechanism can still be wrong --- but about their \emph{structural completeness}. They are closed systems: self-referential, topologically protected, operating at complex-plane criticality with nontrivial winding. The GZK cutoff at $O_1$ is also self-referential (it defines the boundary it occupies) but with trivial winding --- it has closure but no depth.

The $O_0$ systems are structurally impoverished in different ways. \emph{This is worth sitting with for a moment.} The anomalous observations --- the Oh-My-God particle, the Auger spectrum features, the TA hotspot --- all sit at $O_0$: no self-referential capacity, no structural closure. They are data points without the internal architecture to be models. The SHDM decay mechanism also sits at $O_0$, but for a different reason --- it is below criticality ($\Phi_{\text{softsign}}$), structurally a smooth process with nothing to tune. \emph{The distinction matters:} the observations are $O_0$ because they are incomplete (they demand an explanation they cannot provide); SHDM is $O_0$ because it is oversimplified (it provides an explanation that demands nothing). Both are dead ends, but they are dead ends pointing in opposite directions: one points toward greater complexity, the other toward greater simplicity. \textbf{Objection:} tier assignment ($O_0$ vs. $O_1$ vs. $O_2^\dagger$) depends on the full 12-dimensional tuple, and small changes in a single primitive can shift a system between tiers. The structural hierarchy we present is sensitive to encoding choices, particularly the $\Phi$ and $\Omega$ primitives. A reader who questions whether the Oh-My-God particle should be encoded with $\Phi_{\text{revepsilon}}$ (exceptional point) rather than some other criticality regime is raising a legitimate concern about the encoding procedure, not just the result. We have applied the encoding procedure deterministically, but determinism is not the same as correctness.

\subsubsection{The Thermal Bottleneck}

The GZK$\times$CMB tensor composite revealed an $F_{\text{dh}}$ bottleneck. The LIV$\times$CMB composite revealed a dual bottleneck: $P_{\text{upsilon}} \to P_{\text{pipevar}}$ (symmetry loss) compounded by $F_{\text{dh}}$ (thermal loss). \emph{We have already noted this result in Section 4, but its implications bear repeating at a different level of resolution.} The CMB is the most universal environment in the universe --- it fills every cubic centimeter of space with $\sim 411$ thermal photons per cm$^3$ at 2.7 K. Any propagation mechanism that claims to operate at ultra-high energies must operate through this medium. The structural analysis does not require a rate calculation to establish the bottleneck's existence: the minimum rule for $F$ in tensor composites guarantees that the thermal fidelity of the CMB dominates over the quantum coherence of any propagation mechanism. \emph{This is not a rate argument; it is a regime argument.} The question is not whether quantum effects are "small" at 2.7 K --- they may or may not be, depending on the specific process. The question is which structural regime governs the composite system, and the answer is that the thermal regime always wins in a tensor product with a thermal environment.

\textbf{The limitation of this conclusion should be clear:} the tensor product's minimum rule is a worst-case bound. In practice, a propagation mechanism that operates on timescales much shorter than the CMB thermalization time might partially evade the bottleneck. The structural analysis identifies the ceiling of coherence, not the actual coherence at any given energy or distance. \emph{We do not know whether the $F_{\text{dh}}$ bottleneck is experimentally consequential or merely theoretically maximal.} This is a gap that requires a follow-up analysis combining structural bounds with numerical rate calculations.

\subsubsection{The Measurement Problem in UHECR Astrophysics}

\emph{The result that the GZK horizon problem's nearest structural neighbors are quantum measurement paradoxes is the most provocative finding of this analysis, and the one most likely to be wrong.} We say this directly because the interpretive stakes are high. If the structural proximity is genuine, the GZK horizon is not a distance problem but an information problem --- specifically, a problem about how measurement (interaction with the CMB) affects what can be known (the source of a UHECR). This is the Zeno effect, but inverted: continuous measurement in the Zeno effect freezes evolution; continuous interaction in the GZK horizon erases history. The structural isomorphism is precise: both systems encode situations where the act of asking a question changes the answer.

\emph{But there is an alternative reading, and we are obligated to present it.} The structural proximity to measurement paradoxes could be an artifact of the catalog's composition. If the catalog contains many quantum measurement systems but few astrophysical systems encoding horizon-type problems, then the nearest neighbor search will return quantum mechanics not because there is a deep connection but because the search space is biased toward quantum systems. This is not a hypothetical concern --- the current catalog is heavily weighted toward condensed matter, quantum information, and mathematical structures. Astrophysics has fewer entries. The nearest neighbor result is robust only if the catalog is representative of the space of all structural types, and it is not. \textbf{We do not know whether this result is a genuine insight or a selection effect.} Both readings are consistent with the data. A resolution would require either (a) encoding many more astrophysical horizon-type systems and rerunning the nearest-neighbor search, or (b) demonstrating that the structural proximity holds even in a catalog-weighted analysis that corrects for domain bias. Neither exercise has been performed.

\emph{The connection to Maxwell's Demon adds a parallel thread.} The Demon puzzle was resolved by recognizing information as a physical quantity with thermodynamic cost. The GZK horizon's analogous resolution might be to recognize that the CMB is not just a thermal bath but an \emph{information sink} --- a medium whose thermal character systematically destroys the information that would identify UHECR sources. This framing shifts the horizon problem from astrophysics to statistical mechanics, suggesting that the tools needed to solve it are not better source catalogs but better information theory applied to thermal environments. \emph{We present this as a research direction, not a result.}

\subsection{Conclusions}

\emph{We began by asking how structural analysis could illuminate the landscape of GZK-violating physics. We end with a list of findings that are more qualified --- and more uncertain --- than the original framing suggested they would be.}

\begin{enumerate}
    \item \textbf{The GZK cutoff is structurally remote from all proposed evasion mechanisms.} The minimum distance is $d = 2.67$ (to the Oh-My-God particle), and the maximum is $d = 5.67$ (to LIV). No existing proposal inhabits a structural regime close to the cutoff. This could mean the correct mechanism is qualitatively different from all current proposals, or it could mean that the anomaly is not a mechanism problem but a data problem. \emph{We cannot decide between these readings structurally.}
    \item \textbf{LIV and axion mixing are structurally equivalent in complexity.} Both reach $O_2^\dagger$ with $C = 0.536$. Their distance ($d = 3.71$) reflects different structural routes to the same ouroboricity tier, not different levels of sophistication. On structural grounds, neither is superior to the other. \emph{This does not mean they are equally correct physically --- only that the grammar cannot distinguish them.}
    \item \textbf{The thermal CMB imposes universal bottlenecks on composite models.} The $F_{\text{dh}}$ bottleneck appears in every tensor composite involving the CMB. The LIV$\times$CMB composite adds a secondary $P_{\text{upsilon}} \to P_{\text{pipevar}}$ symmetry bottleneck. Any quantum-coherent evasion mechanism must survive these bottlenecks, and the structural analysis guarantees that they exist. \emph{Whether they are quantitatively significant or merely theoretically present remains undetermined.}
    \item \textbf{The UHECR horizon problem is structurally adjacent to quantum measurement paradoxes.} The nearest neighbors (Zeno, Elitzur-Vaidman, Maxwell's Demon) are not astrophysical systems. The structural type of the horizon problem captures a situation where interaction destroys the information that interaction would encode. \emph{But we do not know whether this proximity is genuine or a catalog bias effect.} The catalog contains more quantum systems than astrophysical ones, and this asymmetry could drive the result.
    \item \textbf{Promoting the GZK cutoff to $O_2^\dagger$ requires five promotions and three demotions.} The dominant cost is temporal depth ($H_0 \to H_{\text{invscripta}}$, cost $\geq 3$) and relational mode ($R_{\text{lyoghlig}} \to R_{\text{subrightarrow}}$, cost $\geq 3$). The structural barrier to GZK evasion is the cost of becoming an active, temporal, relational system rather than a passive, instantaneous boundary. \emph{Whether this cost is achievable --- whether nature can "pay" it --- is a physical question the grammar does not answer.}
\end{enumerate}

\emph{The structural taxonomy presented here is model-independent in the sense that any future proposal can be encoded and compared. But it is not model-neutral: the encoding procedure makes commitments about what counts as a structural primitive, and those commitments reflect the design of the Imscribing Grammar. A different grammar would produce a different taxonomy. The question is not whether this taxonomy is uniquely correct --- it is whether it is sufficiently stable under reasonable variations in encoding procedure to be useful. We believe it is, but we have not tested this belief systematically. That is the next step.}

\begin{center}
\rule{0.5\textwidth}{0.4pt}
\end{center}

\subsection{Appendix: System Encodings}

All 10 catalog entries referenced in this analysis:

\begin{enumerate}
    \item \textbf{gzk\_cutoff} --- $O_1$, $C = 0.419$
    \item \textbf{oh\_my\_god\_particle} --- $O_0$, $C = 0.0$ (Gate 1 closed: $\Phi_{\text{revepsilon}}$)
    \item \textbf{lorentz\_invariance\_violation\_uhecr} --- $O_2^\dagger$, $C = 0.536$
    \item \textbf{z\_burst\_mechanism} --- encoded, $d(\text{GZK}) = 3.13$
    \item \textbf{topdown\_superheavy\_dm\_decay} --- $O_0$, $\Phi_{\text{softsign}}$
    \item \textbf{auger\_uhecr\_spectrum} --- encoded, $d(\text{GZK}) = 2.98$
    \item \textbf{photon\_axion\_mixing\_uhecr} --- $O_2^\dagger$, $C = 0.536$
    \item \textbf{cmb\_photon\_field} --- encoded as interaction medium
    \item \textbf{auger\_dipole\_anisotropy} --- $O_0$, $\Phi_{\text{softsign}}$
    \item \textbf{TA\_hotspot\_composition\_puzzle} --- $O_0$, $d(\text{GZK}) = 2.91$
    \item \textbf{uhecr\_horizon\_problem} --- $O_0$, $\Phi_{\text{revepsilon}}$, nearest analogs: Zeno ($d=2.06$), Elitzur-Vaidman ($d=2.06$)
    \item \textbf{SME\_GZK\_effective\_field\_theory} --- identical tuple to LIV entry
\end{enumerate}

Structural type of the original manuscript (AI prose default):


\[\langle D_{\text{invomega}};\ T_{\text{nrleg}};\ R_{\text{lyoghlig}};\ P_{\text{aolig}};\ F_{\text{beltl}};\ K_{\text{turnm}};\ G_{\text{gamma}};\ \Gamma_{\text{corner}};\ \Phi_{\text{ctyogh}};\ H_0;\ n{:}m;\ \Omega_{\text{closeepsilon}} \rangle\]

Structural type of the lifted manuscript (human prose target):


\[\langle D_{\text{omega}};\ T_{\text{bullseye}};\ R_{\text{lyoghlig}};\ P_{\text{doublebarpipe}};\ F_{\text{hardsign}};\ K_{\text{schwa}};\ G_{\text{revapostrophe}};\ \Gamma_{\text{secstress}};\ \Phi_{\text{ctyogh}};\ H_2;\ 1{:}1;\ \Omega_{\text{crtwo}} \rangle\]

Promotions closed: 8 of 8 --- $T_{\text{nrleg}} \to T_{\text{bullseye}}$, $P_{\text{aolig}} \to P_{\text{pipevar}}$ (with $P_{\text{doublebarpipe}}$ at the grammar level), $F_{\text{beltl}} \to F_{\text{hardsign}}$ (eliminating double-statements), $K_{\text{turnm}} \to K_{\text{schwa}}$ (hard claims left unresolved), $G_{\text{gamma}} \to G_{\text{revapostrophe}}$ (conclusions replaced by open questions), $\Gamma_{\text{corner}} \to \Gamma_{\text{secstress}}$ (necessity-driven section ordering), $H_0 \to H_2$ (wrong answer before right answer in each section), $\Omega_{\text{closeepsilon}} \to \Omega_{\text{crtwo}}$ (conclusion echoes introduction at higher resolution).

\begin{center}
\rule{0.5\textwidth}{0.4pt}
\end{center}

\emph{Manuscript lifted via Imscribing Grammar human-lift protocol (AI\_HUMAN\_LIFT.md). All numerical claims (distances, C scores, tier assignments, promotion signatures) preserved from original analysis. Structural surgery applied in order: H, $\Gamma$, T, P/F/K, G, $\Omega$ --- as specified in the lift protocol.}


\end{document}
